\section{Technology Stack}
\label{sec:implementation:tech-stack}

The system is implemented on a modern, cross-platform technology stack optimized for real-time control, maintainability, and extensibility. At its core, Qt with C++ provides both the backend logic and the graphical interface, leveraging its signal-slot architecture to ensure asynchronous communication and deterministic event handling. This is complemented by QML, which extends the framework with a flexible and animated simulation environment, allowing operators and developers to visualize interlocking behavior dynamically. Persistent data management is handled through PostgreSQL, which maintains normalized representations of both static infrastructure and runtime interlocking states, ensuring consistency and traceability. Build processes are coordinated through CMake, enabling reproducible configurations and streamlined compilation across platforms. Together, these technologies form a cohesive ecosystem that tightly integrates logic processing, data persistence, and visualization—an essential combination for an interlocking system where safety, usability, and adaptability must coexist seamlessly.
